\documentclass[12 pt, openany]{book}
\usepackage{preamble}

% --- GLOSSARY --- %
\usepackage{datatool}[=v2.32] % pervent error with babel package
\usepackage[
    %nonumberlist,      % remove numbers reference in glossaries
    counter=page,       % set the counter where glossary entry are defined to point at where the term was referenced
    nogroupskip,        % define the same space for each glossarie entry
    section=subsection, % Define glossary outline level
]{glossaries}
\usepackage[
    toc,                % include glossary in the table of content
    acronym,            % to use the acronym
    abbreviations,      % to use abbreviations
    postdot,            % adds a dot after the explanation of the glossary entry,
    noredefwarn,        %suppress the warning that occurs when the environment and \printglossary are redefined while glossaries is loading
    section=subsection, % Define glossary outline level
    numberedsection,     % Set the numbering as an usual subsection of the
]{glossaries-extra} % improves glossaries package feature

\usepackage{glossaries}

% Redefine how the entry name is printed for the 'unsrt' style
\renewcommand*{\glossentryname}[1]{%
  \Glsentrytext{#1}%
}

\makenoidxglossaries
\loadglsentries{glossary.tex}

% --- MAIN DOCUMENT --- %
\begin{document}
    \frontmatter
    \import{./}{title}

    \setcounter{tocdepth}{2}
    \tableofcontents

    \mainmatter
    
    \setcounter{fnm}{\value{footnote}}

    % Section

    \subfile{section/introduction}

    \subfile{section/architectural_design}

    \subfile{section/user_interface_design}

    \subfile{section/requirements_traceability}

    \subfile{section/implementation_integration_testplan}

    \subfile{section/effort_spent}

    \backmatter

    %\printunsrtglossary[type=main, title=Glossary, toctitle=Glossary] % Print the glossary entries
    %\printunsrtglossary[type=\acronymtype] % Print acronym entries

    \nocite{*}
    
    \printbibliography[heading=bibintoc] % Print all the bibliography and adds it to the list of contents
    %\printbibliography[type=article, title={Articles Reference}] % Prints the bibliography references related to Articles
    %\printbibliography[type=book, title={Books Reference}] % Prints the bibliography references related to Books
    %\printbibliography[type=manual, title={Manual Reference}] % Prints the bibliography references related to Manual
    %\printbibliography[type=online, title={Online Reference}] % Prints the bibliography references related to online soruces
\end{document}